\section{设计优先的Vibe Coding方法}
\label{sec:method}

我们提出\textbf{设计优先的Vibe Coding}，这是一套面向非专业用户的五步方法论，通过与AI协作产出高质量的网页应用。
其核心原则是：\textit{将用户的角色从设计者转变为设计评审者}——AI提出设计方案，用户进行判断，最终收敛的规范成为实现的蓝图。

\subsection{第一步：确立设计哲学}

用户和AI不应急于生成代码，而是首先通过深入对话确定项目的\textit{设计哲学}——一个将约束所有后续决策的抽象框架。

在我们的实现中，这段对话产生了一个核心矛盾：\textbf{温润 $\times$ 锐气}。
温润通过材质选择来体现（奶油色纸张质感底色、衬线字体、充裕的留白），而锐气通过工艺精度来体现（精确编排的动画、横向翻阅交互、分层的微交互）。
AI将其总结为设计约束：所有动画时长控制在0.3-0.8秒以内、采用快入慢出的缓动曲线、每个动效都服务于明确的目的。

\subsection{第二步：定义信息架构}

设计哲学确立后，用户和AI共同规划页面的\textit{信息架构}——存在哪些区域、以什么顺序排列、以及它们为访客创造了怎样的情感节奏。

我们的实现定义了六个按序排列的区域：（1）首屏——名字与态度宣言，配合氛围粒子；（2）关于——头像、自我介绍与可交互的兴趣标签；（3）思考——横向翻阅文章卡片（"招牌交互"）；（4）手记——瀑布流布局与分类筛选；（5）联系——社交图标与悬停微交互；（6）页脚——分层彩蛋系统。

\subsection{第三步：逐区域细化}

针对每个区域，用户和AI对以下内容进行具体约定：视觉元素（出现什么）、交互规格（点击/悬停/滚动时发生什么）、动画细节（时长、缓动、交错）。

正是在这一步，AI的设计能力最直接地弥补了用户的设计局限。
例如，当用户提出"想要一种像翻书脊一样的横向滑动"时，AI提出了基于ScrollTrigger的pin+scrub实现方案，分析了哪个区域最适合这一交互，并细化了卡片聚焦效果、导航圆点以及移动端降级策略。

\subsection{第四步：编写设计文档}

将前三步的所有决策编写为一份结构化的Markdown文档（我们的案例中约700行）。
该文档在第五步中充当\textit{唯一真相来源}。
关键的是，它既可供人类用户阅读（用于验证设计意图），也可供AI阅读（在实现过程中作为参考约束）。

\subsection{第五步：依据规范实现}

只有当设计文档完成后，代码生成才开始。
AI按照文档逐区域实现，将设计文档作为约束参考。
当实现偏离规范时，文档为纠正提供了客观参照。

这种设计决策与编码实现的分离，是设计优先策略区别于增量构建的关键所在——\textit{设计决策在代码之前做出，而非在代码中摸索。}

\subsection{迭代循环}

第五步产出一个可运行的版本后，用户在浏览器中进行实际测试。
设计文档与实际体验之间的差异——这是不可避免的，因为某些品质（如视觉权重、动画手感）无法在文字中完全描述——会通过反馈给AI进行针对性修正。
这一循环与增量构建的"添加模块、调整、重复"模式有本质区别：修正\textit{以设计文档为锚点}，而非无参照的任意调整。
文档将修复空间约束在可控范围内，防止风格漂移。
